% © 2025 Ing. David Jaroš — CC BY-NC-ND 4.0
%
% This work is licensed under a Creative Commons Attribution-NonCommercial-NoDerivatives
% 4.0 International License (CC BY-NC-ND 4.0).
%
% License History: Earlier drafts (up to v0.3) were released under CC BY 4.0.
% From v0.4 onward, all material is released under CC BY-NC-ND 4.0 to protect
% the integrity of the theoretical work during ongoing academic development.
%
% See LICENSE.md for full license text.

\documentclass[11pt,a4paper]{article}
\usepackage{amsmath, amssymb, amsthm}
\usepackage{hyperref}
\usepackage[a4paper, margin=2.5cm]{geometry}
\usepackage{tcolorbox}
\usepackage{longtable}
\usepackage{booktabs}

\newtheorem{theorem}{Theorem}[section]
\newtheorem{proposition}[theorem]{Proposition}
\newtheorem{lemma}[theorem]{Lemma}
\newtheorem{corollary}[theorem]{Corollary}
\theoremstyle{definition}
\newtheorem{definition}[theorem]{Definition}
\theoremstyle{remark}
\newtheorem{remark}[theorem]{Remark}

\title{\textbf{Experimental Predictions of UBT}\\
\large Unified Biquaternion Theory — Canonical Appendices Series}
\author{Ing.\ David Jaroš}
\date{March 2026}

\begin{document}
\maketitle
\tableofcontents
\bigskip

\begin{tcolorbox}[colback=blue!4!white,colframe=blue!50!black,title=Purpose]
This appendix extracts concrete predictions from UBT that are in principle
measurable or that distinguish UBT from Standard Model / GR.
It classifies predictions by their derivation status and identifies the
most promising observational tests.

\medskip
\textbf{Sources}:
\texttt{canonical/qed\_phi\_const/},
\texttt{canonical/qm\_emergence/},
\texttt{docs/reports/hecke\_lepton/},
\texttt{DERIVATION\_INDEX.md}

\medskip
\begin{tabular}{lll}
\hline
Category & Prediction count & Section \\
\hline
Electrodynamics & 4 & \S\ref{sec:qed} \\
Fine structure constant & 3 & \S\ref{sec:alpha_pred} \\
Particle spectrum & 4 & \S\ref{sec:spectrum} \\
Cosmological & 3 & \S\ref{sec:cosmo} \\
Imaginary-sector deviations & 3 & \S\ref{sec:imaginary} \\
\hline
\end{tabular}
\end{tcolorbox}

\begin{tcolorbox}[colback=yellow!5!white,colframe=orange!60!black,title=Prediction classification]
Predictions are labelled:
\begin{itemize}
  \item \textbf{[DERIVED]} — follows from UBT without free parameters at current layer
  \item \textbf{[SEMI-EMP]} — structural derivation with $\geq 1$ calibrated parameter
  \item \textbf{[CONJECTURE]} — proposed; not yet derived
  \item \textbf{[CONSISTENT]} — same as SM/GR; UBT does not modify this observable
\end{itemize}
\end{tcolorbox}

% -----------------------------------------------------------------
\section{Electrodynamics Predictions}\label{sec:qed}
% -----------------------------------------------------------------

\subsection{No Photon Mass}

\begin{theorem}[$\mathrm{U}(1)_{\mathrm{EM}}$ protection]\label{thm:no_mass}
The action $S[\Theta]$ preserves exact $\mathrm{U}(1)_{\mathrm{EM}}$ symmetry
at all loop orders.  Therefore, the photon mass $m_\gamma = 0$ exactly.
\end{theorem}

\noindent\textbf{Status}: \textbf{[DERIVED]} [L0] \\
\noindent\textbf{Source}: \texttt{canonical/qed\_phi\_const/step1\_u1\_protection.tex} \\
\noindent\textbf{Experimental bound}: $m_\gamma < 10^{-18}$ eV (consistent).

\subsection{Schwinger Anomalous Magnetic Moment}

\begin{proposition}[Schwinger term at $\varphi = \mathrm{const}$]
\label{prop:schwinger}
In the QED limit ($\varphi = \mathrm{const}$, real sector), the anomalous
magnetic moment of the electron is:
\begin{equation}\label{eq:schwinger}
  a_e = \frac{\alpha}{2\pi} + O(\alpha^2)
  \approx 1.161 \times 10^{-3}.
\end{equation}
\end{proposition}

\noindent\textbf{Status}: \textbf{[DERIVED]} [L1] (Schwinger term proved) \\
\noindent\textbf{Source}: \texttt{canonical/qed\_phi\_const/step4\_schwinger\_term.tex} \\
\noindent\textbf{Experiment}: $a_e^{\mathrm{exp}} = 1.15965\ldots \times 10^{-3}$ (consistent).

\subsection{QED Running Coupling ($\beta$-function)}

\begin{proposition}[$\beta$-function at one loop]
\label{prop:beta}
The one-loop $\beta$-function for the electromagnetic coupling is:
\begin{equation}\label{eq:beta}
  \mu\frac{\mathrm{d}\alpha}{\mathrm{d}\mu}
  = \frac{2\alpha^2}{3\pi} + O(\alpha^3).
\end{equation}
\end{proposition}

\noindent\textbf{Status}: \textbf{[DERIVED]} [L1] \\
\noindent\textbf{Source}: \texttt{canonical/qed\_phi\_const/step3\_beta\_function.tex} \\
\noindent\textbf{Experiment}: Standard one-loop QED $\beta$-function verified to $< 0.1\%$.

\subsection{Lamb Shift}

The Lamb shift in hydrogen ($2S_{1/2} - 2P_{1/2}$ splitting) receives a
leading contribution from vacuum polarisation at one loop:
\begin{equation}
  \Delta E_{\mathrm{Lamb}} \approx 1058\;\mathrm{MHz}
  \quad (\text{QED prediction}).
\end{equation}

\noindent\textbf{Status}: \textbf{[CONSISTENT]} (UBT QED limit reproduces standard QED) \\
\noindent\textbf{Source}: \texttt{canonical/qed\_phi\_const/step5\_lamb\_shift.tex} (sketch).

% -----------------------------------------------------------------
\section{Fine Structure Constant Predictions}\label{sec:alpha_pred}
% -----------------------------------------------------------------

\subsection{Bare Value $\alpha^{-1} = 137$}

\begin{proposition}
From the prime attractor of $V_{\mathrm{eff}}(n)$:
\begin{equation}
  \alpha^{-1}_{\mathrm{bare}} = n^* = 137 \quad (\text{prime}).
\end{equation}
\end{proposition}

\noindent\textbf{Status}: \textbf{[SEMI-EMP]} (B-coefficient partially derived) \\
\noindent\textbf{Source}: \texttt{docs/STATUS\_ALPHA.md}, \texttt{appendix\_alpha\_geometry.tex}.

\subsection{Full Value $\alpha^{-1} = 137.036$}

The full inverse fine-structure constant including one-loop QED correction:
\begin{equation}
  \alpha^{-1} = 137 + \delta, \qquad \delta \approx 0.036.
\end{equation}

\noindent\textbf{Status}: \textbf{[SEMI-EMP]} (correction $\delta$ from QED running; $R$-factor open) \\
\noindent\textbf{Source}: \texttt{docs/STATUS\_ALPHA.md \S6}.

\subsection{Non-circularity Test}

Different values of $N_{\mathrm{eff}}$ (i.e., different SM particle content)
give different prime attractors $n^*$:
\begin{equation}
  N_{\mathrm{eff}} = 1 \;\Rightarrow\; n^* \neq 137.
\end{equation}
This test verifies that the result is not circular.

\noindent\textbf{Status}: \textbf{[DERIVED]} \\
\noindent\textbf{Source}: \texttt{experiments/validation/validate\_B\_coefficient.py}.

% -----------------------------------------------------------------
\section{Particle Spectrum Relations}\label{sec:spectrum}
% -----------------------------------------------------------------

\subsection{SU(3) Color Sector: Confinement}

\begin{tcolorbox}[colback=green!4!white,colframe=green!50!black,title=Algebraic prediction]
Free quarks (colour non-singlets) are algebraically inadmissible in UBT:
$\mathcal{H}_{\mathrm{phys}} = \mathrm{Im}(\Pi_{\mathrm{color}})$ in
$\mathbb{C}^2\otimes\mathbb{C}^2\otimes\mathbb{C}^2$.
All verified hadrons (baryons, mesons, tetraquarks, pentaquarks) satisfy
$\langle C_2 \rangle = 0$.
\end{tcolorbox}

\noindent\textbf{Status}: \textbf{[DERIVED]} (algebraic confinement) / Clay Millennium (mass gap) \\
\noindent\textbf{Source}: \texttt{ARCHIVE/archive\_legacy/consolidation\_project/confinement/}.

\subsection{Lepton Mass Ratios}

The Hecke structure at prime $p = 137$ gives strong numerical support for:
\begin{equation}
  \frac{m_\mu}{m_e} \approx 207,
  \qquad
  \frac{m_\tau}{m_\mu} \approx 16.8.
\end{equation}

\noindent\textbf{Status}: \textbf{[CONJECTURE]} (numerical support, no analytic derivation) \\
\noindent\textbf{Source}: \texttt{docs/reports/hecke\_lepton/}.

\subsection{Mirror Sector at $p^{**} = 139$}

The twin prime $p^{**} = 139$ gives a mirror-sector prediction for a second set
of mass ratios (Set B):
$(m_1/m_e, m_2/m_e, m_3/m_e) \approx (195.2, 50.4, 54.6)$.

\noindent\textbf{Status}: \textbf{[CONJECTURE]} (numerical observation; distinct from SM) \\
\noindent\textbf{Source}: \texttt{docs/reports/hecke\_lepton/mirror\_world\_139.md}.

\subsection{Three Generations from Quaternionic Structure}

The three imaginary quaternion units $\{I,J,K\}$ of $\mathbb{H}$ correspond to
exactly three fermion families (no more, no fewer).

\noindent\textbf{Status}: \textbf{[DERIVED]} (algebraic [L1]); \\
\noindent\textbf{Source}: \texttt{research\_tracks/research/generation\_structure.tex}.

% -----------------------------------------------------------------
\section{Cosmological Predictions}\label{sec:cosmo}
% -----------------------------------------------------------------

\subsection{Hubble Tension (Metric Latency)}

The imaginary sector $h_{\mu\nu} = \mathrm{Im}(\mathcal{G}_{\mu\nu})$ introduces
a \emph{metric latency}: the effective Hubble constant inferred from late-time
observations differs from the early-universe value by a correction
$\delta H / H \sim |h_{\mu\nu}|/|g_{\mu\nu}|$.

\noindent\textbf{Status}: \textbf{[CONJECTURE]} (speculative; qualitatively consistent with $H_0$ tension) \\
\noindent\textbf{Source}: \texttt{speculative\_extensions/appendices/appendix\_HT\_hubble\_tension\_metric\_latency.tex}.

\subsection{Dark Sector from $p$-adic Extensions}

$p$-adic extensions of the $\psi$-winding structure give stable, non-interacting
modes that do not couple to $\mathrm{U}(1)_{\mathrm{EM}}$ — a candidate for
dark matter.

\noindent\textbf{Status}: \textbf{[CONJECTURE]} (speculative) \\
\noindent\textbf{Source}: \texttt{research\_tracks/p\_universes/}.

\subsection{Cosmological GR Limit}

All standard cosmological solutions of GR (FRW metric, Schwarzschild metric,
gravitational waves) are exactly reproduced on the admissible sector
$\mathcal{A}_{\mathrm{UBT}}$.

\noindent\textbf{Status}: \textbf{[CONSISTENT]} [L1] \\
\noindent\textbf{Source}: \texttt{canonical/gr\_limit/GR\_limit\_of\_UBT.tex}.

% -----------------------------------------------------------------
\section{Imaginary-Sector Deviations}\label{sec:imaginary}
% -----------------------------------------------------------------

These are novel UBT predictions beyond Standard Model + GR, arising from the
imaginary sector $h_{\mu\nu}$ and the phase curvature.

\subsection{Phase Curvature Signatures}

Corrections of order $|h_{\mu\nu}/g_{\mu\nu}| \sim |\psi|^2/R_\psi^2$ to
gravitational wave phase shifts and light-bending angles.

\noindent\textbf{Status}: \textbf{[CONJECTURE]} (order of magnitude estimate only) \\
\noindent\textbf{Source}: \texttt{canonical/geometry/biquaternion\_metric.tex}.

\subsection{$\phi$-Universe Parameter}

The scalar phase $\phi$ of $\Theta$ is physical (not pure gauge) for two-mode
winding vacua, introducing a moduli parameter $\phi \in [0,2\pi)$ that could
in principle be measured through interference experiments.

\noindent\textbf{Status}: \textbf{[DERIVED]} (existence of physical $\phi$ proved) \\
\noindent\textbf{Source}: \texttt{docs/PHI\_UNIVERSE\_PARAMETER.md \S4a}.

\subsection{Closed Timelike Curves (Speculative)}

Certain non-generic solutions of the master equation may admit closed timelike
curves; this is a speculative extension separated from canonical UBT.

\noindent\textbf{Status}: \textbf{[SPECULATIVE]} \\
\noindent\textbf{Source}: \texttt{speculative\_extensions/appendices/appendix\_J\_rotating\_spacetime\_ctc.tex}.

% -----------------------------------------------------------------
\section{Priority Predictions Summary}\label{sec:priority}
% -----------------------------------------------------------------

The following predictions are most likely to distinguish UBT from the
Standard Model in near-term experiments:

\begin{center}
\begin{longtable}{llll}
\toprule
Prediction & Value & Status & Distinguishes UBT? \\
\midrule
No photon mass & $m_\gamma = 0$ & DERIVED & No (consistent) \\
Schwinger term & $a_e = \alpha/(2\pi)$ & DERIVED & No (consistent) \\
$\beta$-function at 1-loop & standard QED & DERIVED & No (consistent) \\
$\alpha^{-1}_{\mathrm{bare}} = 137$ & 137 & SEMI-EMP & Partially (origin) \\
Three generations & 3 families & DERIVED & Partially \\
Algebraic confinement & no free quarks & DERIVED & No (consistent) \\
Mirror sector at $p=139$ & mass spectrum B & CONJECTURE & \textbf{Yes} \\
$\phi$-parameter physical & $\phi \neq$ pure gauge & DERIVED & \textbf{Yes} \\
Phase curvature & $\delta\phi_{\mathrm{grav}}\sim h/g$ & CONJECTURE & \textbf{Yes} \\
Hubble tension correction & $\delta H/H \sim h/g$ & CONJECTURE & \textbf{Yes} \\
\bottomrule
\end{longtable}
\end{center}

\end{document}
